\section{Read and Write Reorder Buffers}
\label{sec:rob}
\subsection{Overview}

The ROB is a generic transaction-tracking core, instantiated independently for the read path and the write path. Each instance owns a complete, separate copy of every structure described in \S\ref{sec:generic} --- its own 256-entry Free List, its own 256-entry Metadata RAM, and its own bank of 64 Per-AXID Slot FIFOs (depth 4 each). Nothing is shared between the two instances. They differ only in what a completed transaction produces: the read instance additionally holds a Data Buffer association and drives a Read Serializer; the write instance holds an aggregate response-status field and a 64\,B line footprint consumed by read admission (\S\ref{sec:samelinehaz}), and drives \texttt{BRESP} generation directly, since a write completion carries no payload.

Each transaction is identified end-to-end by a single \texttt{GLOBAL\_TAG}, described in \S\ref{sec:tag}.  \texttt{ROB\_INDEX} is the Metadata RAM address end to end.

A transaction becoming \emph{complete} (all its packets accounted for in \texttt{completion\_bitmap}) is not the same event as a transaction being \emph{freed}. A completed transaction remains allocated while its ordered AXI response is being emitted and accepted, and (read only) while its contiguous RD\_SRAM range is still being drained. \S\ref{sec:retire} makes this lifecycle explicit.

\subsection*{Lifecycle Summary}
The full life of a ROB entry, in order, referenced to the sections that detail each stage:
\begin{enumerate}[noitemsep]
  \item AXI request accepted (upstream of this file) --- \S\ref{sec:alloc}
  \item \texttt{ROB\_INDEX} allocated from the Free List --- \S\ref{sec:freelist}, \S\ref{sec:alloc}
  \item Metadata RAM initialized (\texttt{exp\_pkts}, \texttt{completion\_bitmap = 0}, (read) \texttt{rd\_dbuf\_base}) --- \S\ref{sec:alloc}
  \item Entry becomes visible in its AXID's Slot FIFO --- \S\ref{sec:fifo}, \S\ref{sec:alloc}
  \item Packets emitted, \texttt{GLOBAL\_TAG} attached --- \S\ref{sec:issue}
  \item Completions decoded, \texttt{completion\_bitmap} updated --- \S\ref{sec:compin}, \S\ref{sec:completion}
  \item Transaction reaches completion; becomes eligible to retire --- \S\ref{sec:retire}
  \item Ordered AXI response emitted and accepted (read: beat by beat; write: single \texttt{BRESP}) --- \S\ref{sec:read}, \S\ref{sec:write}
  \item Slot FIFO popped, \texttt{ROB\_INDEX} freed, (read) RD\_SRAM range already released --- \S\ref{sec:retire}, \S\ref{sec:release}
\end{enumerate}

\subsection*{Architecture principles}
\begin{itemize}[noitemsep]
  \item Completion directly supplies \texttt{ROB\_INDEX} and \texttt{PACKET\_NUMBER}; RAM access is a bit-slice decode.
  \item Per-AXID ordering is carried entirely by the Per-AXID Slot FIFOs (\S\ref{sec:fifo}).
  \item Packet position is the low 6 bits of \texttt{GLOBAL\_TAG} (\texttt{PACKET\_NUMBER}).
  \item Read request-issue backpressure is a function of the pooled Read Data Buffer's own 128-entry watermark, a separate structure outside this file's scope (see \S\ref{sec:backpressure}). No watermark applies to Metadata RAM.
\end{itemize}

\section{Global Tag}
\label{sec:tag}

\begin{center}
\begin{tabular}{ll}
\toprule
\textbf{Field} & \textbf{Width} \\
\midrule
\texttt{ROB\_INDEX}     & 8 bits (addresses one of 256 Metadata RAM entries) \\
\texttt{PACKET\_NUMBER} & 6 bits (64-B MC packet index, \texttt{0..63}) \\
\midrule
\texttt{GLOBAL\_TAG}    & 14 bits, \texttt{\{ROB\_INDEX[7:0], PACKET\_NUMBER[5:0]\}} \\
\bottomrule
\end{tabular}
\end{center}

\textbf{Generation.} At allocation, the Allocator (\S\ref{sec:freelist}) issues \texttt{ROB\_INDEX}. At packet emission, the Burst Splitter supplies \texttt{PACKET\_NUMBER} (\texttt{0..exp\_pkts-1}) for each 64-B MC packet of the transaction. Packet Formation concatenates the two into \texttt{GLOBAL\_TAG}, which is attached to every packet sent to MC Core.

\textbf{Carry.} MC Core is not required to interpret \texttt{GLOBAL\_TAG} --- it echoes the value back unmodified, \emph{exactly one echo per accepted 64-B request packet}, as required by \S\ref{sec:cifmc-req} (item 1). The \texttt{status[3:0]} to AXI \texttt{RESP} mapping is frozen in \S\ref{sec:cifmc}.

\textbf{Decode.} On completion, Completion Input (\S\ref{sec:compin}) splits the returned \texttt{GLOBAL\_TAG} directly: bits \texttt{[13:6]} address Metadata RAM (\texttt{ROB\_INDEX}); bits \texttt{[5:0]} select the packet position within that entry (\texttt{PACKET\_NUMBER}). Both are pure bit-slices --- no lookup, no associative match.

\section{Generic Architecture}
\label{sec:generic}

The blocks in this section are instantiated once per ROB instance (i.e.\ once for read, once for write), with identical logic and identical sizing in both --- see Overview.

\subsection{Allocator / Free List}
\label{sec:freelist}
\begin{itemize}[noitemsep]
  \item 256-bit free/occupied bitmask, one bit per \texttt{ROB\_INDEX}. \texttt{1 = free}, \texttt{0 = occupied}.
  \item Allocate: priority encoder selects lowest free bit, clears it, and outputs \texttt{ROB\_INDEX[7:0]}.
  \item Free: Retirement Logic sets the bit at \texttt{ROB\_INDEX} via \texttt{free\_req}, only once a slot has fully retired (\S\ref{sec:retire}) --- never merely on completion.
  \item \texttt{FREE\_LIST\_EMPTY} (all bits 0) is a raw pool-wide allocation-backpressure signal. With \texttt{ROB\_DEPTH = 256} and 64 $\times$ depth-4 Slot FIFO capacity, it means every transaction slot is allocated and all FIFO capacity is collectively exhausted. A single full Slot FIFO does not imply this condition. See \S\ref{sec:params}.
\end{itemize}

\subsection{Metadata RAM}
\label{sec:ram}
\begin{itemize}[noitemsep]
  \item 256-entry RAM, directly addressed by \texttt{ROB\_INDEX} (no CAM, no translation).
  \item Depth is fixed at \texttt{ROB\_DEPTH = 256}; no watermark or occupancy threshold applies to this RAM (see Overview). Allocation is gated by Free List exhaustion or per-AXID Slot FIFO fullness (\S\ref{sec:backpressure}); these are distinct conditions.
  \item Entry contents are defined explicitly in \S\ref{sec:entry}.
  \item Write port A (allocation): Allocator writes \texttt{exp\_pkts}, clears \texttt{completion\_bitmap}, (read instance) writes \texttt{rd\_dbuf\_base}, and (write instance) writes the \texttt{wr\_line\_base} and \texttt{wr\_line\_count} footprint (\S\ref{sec:samelinehaz}) at \texttt{ROB\_INDEX} on allocation, before the entry is exposed to retirement --- see \S\ref{sec:alloc} for the exact ordering.
  \item Write port B (completion): Completion Input sets one \texttt{completion\_bitmap} bit (both instances) and, for the write instance, folds \texttt{status} into \texttt{resp\_status}, at \texttt{ROB\_INDEX}. \texttt{rd\_dbuf\_base} is not written on completion --- it is consumed to place the returned 64-B data unit at \texttt{rd\_dbuf\_base + PACKET\_NUMBER}.
  \item Read port: Retirement Logic reads the entry at the head-of-FIFO \texttt{ROB\_INDEX} for each AXID to evaluate completion.
  \item No explicit \texttt{valid} bit. A Metadata RAM entry is only meaningful while its \texttt{ROB\_INDEX} is occupied (cleared in the Free List) and present in a Per-AXID Slot FIFO; both conditions are tracked externally to the RAM.
\end{itemize}

\subsection{Per-AXID Slot FIFO}
\label{sec:fifo}
\begin{itemize}[noitemsep]
  \item 64 independent FIFOs, one per AXID, depth 4 each --- frozen sizing, identical for both instances. Each entry stores \texttt{ROB\_INDEX[7:0]}.
  \item Push: only after the corresponding Metadata RAM entry has been initialized (\texttt{exp\_pkts}, \texttt{completion\_bitmap = 0}, (read) \texttt{rd\_dbuf\_base}) --- see \S\ref{sec:alloc}. An entry is never visible in a Slot FIFO before it is safe to read.
  \item Pop: not on completion. Only when the transaction at the head has fully retired --- i.e.\ its ordered AXI response has been emitted and accepted (\S\ref{sec:retire}).
  \item FIFO order = allocation order = required AXI response order for that AXID, by construction. No separate ordering structure is needed.
  \item FIFO-full for a given AXID is the per-AXID backpressure condition on allocation. Because a completed-but-not-yet-retired entry remains at the head until fully drained, the head slot may stay occupied for several cycles after completion during response emission --- this is expected, not a fault condition.
\end{itemize}

\subsection{Completion Input}
\label{sec:compin}
\begin{itemize}[noitemsep]
  \item Receives the MC completion/response unit together with
      \texttt{GLOBAL\_TAG} and \texttt{status}. The exact response
      transport, payload width, and handshake are defined by the
      MC--CIF interface contract.
  \item Decodes \texttt{GLOBAL\_TAG} into \texttt{ROB\_INDEX} and \texttt{PACKET\_NUMBER} per \S\ref{sec:tag}.
  \item Sets \texttt{completion\_bitmap[PACKET\_NUMBER] = 1} at \texttt{ROB\_INDEX}.
  \item Read instance: writes the returned 64-B data unit to \texttt{rd\_dbuf\_base + PACKET\_NUMBER} in RD\_SRAM (the base was recorded at allocation, not derived here).  The Data Buffer write protocol is defined by the MC/Data-Buffer interface contract --- see \S\ref{sec:read}.
  \item Write instance: folds \texttt{status} into \texttt{resp\_status} (worst-of; precedence \texttt{DECERR} $>$ \texttt{SLVERR} $>$ \texttt{OKAY}, encoding frozen in \S\ref{sec:cifmc}) --- see \S\ref{sec:write}.
  \item Completion acceptance follows the MC--CIF response interface
      contract. The ROB does not impose a separate completion-side
      handshake beyond that contract.
  \item Completion Input only ever marks \emph{completion}. It never pops a Slot FIFO and never frees a \texttt{ROB\_INDEX} --- both are the exclusive responsibility of Retirement Logic (\S\ref{sec:retire}).
\end{itemize}

\subsection{Retirement Logic}
\label{sec:retire}

Completion and retirement are distinct events. Completion (\texttt{completion\_bitmap} full) makes a transaction \emph{eligible} to retire; it does not by itself release anything. Release happens only after the transaction's ordered AXI response has been fully accepted.

\begin{itemize}[noitemsep]
  \item \textbf{Retiring Register} (generic, minimal addition): a single register per instance holding the \texttt{ROB\_INDEX} currently being drained, plus read-instance progress state consisting of a 6-bit
\texttt{PACKET\_NUMBER} packet pointer and any within-beat
progress state required by the MC/Data-Buffer interface contract.
The write instance requires no analogous data-stream progress state. At most one transaction is draining at a time per instance, consistent with the single response-channel path each instance already serves.
  \item Each cycle, while the Retiring Register is idle, Retirement Logic reads the Metadata RAM entry at the head \texttt{ROB\_INDEX} of every AXID's Slot FIFO and checks that \texttt{completion\_bitmap} bits \texttt{[exp\_pkts-1:0]} are all set (for \texttt{exp\_pkts = 64}, the all-ones condition).
  \item When idle, a priority encoder selects the lowest-numbered AXID whose head entry is complete and loads its \texttt{ROB\_INDEX} into the Retiring Register, entering the \emph{retiring} state. \textbf{The Slot FIFO is not popped and \texttt{ROB\_INDEX} is not freed at this point.}
  \item While retiring, instance-specific emission proceeds (\S\ref{sec:read} for read, \S\ref{sec:write} for write) until the transaction's full AXI response has been accepted downstream (\texttt{merge\_ready} for read, \texttt{bresp\_ready} for write --- \S\ref{sec:io}).
  \item Only once the full response has been accepted: Retirement Logic pops that AXID's Slot FIFO entry, returns \texttt{ROB\_INDEX} to the Free List (\texttt{free\_req}), and the Retiring Register returns to idle.
  \item Maximum one transaction draining at a time per instance; draining a single read transaction spans one AXI beat per 64-B MC packet, up to \texttt{exp\_pkts} $\le$ 64 beats, unlike the previous single-cycle-emit assumption.
\end{itemize}

\section{Read-Instance Extension}
\label{sec:read}

State and logic in this section exist only in the read-path ROB instance.

\subsection{Metadata RAM Extension}
\begin{itemize}[noitemsep]
  \item \texttt{rd\_dbuf\_base[RD\_SRAM\_ADDR\_W-1:0]}: base address of a single contiguously-allocated RD\_SRAM range of \texttt{exp\_pkts} 64-B lines. The returned data for packet \texttt{PACKET\_NUMBER} occupies \texttt{rd\_dbuf\_base + PACKET\_NUMBER}. Written once at allocation (\S\ref{sec:alloc}); \emph{not} written by Completion Input, which consumes it to place the response.
  \item \texttt{RD\_SRAM\_ADDR\_W} is \textbf{TBD}: RD\_SRAM depth is not specified in any current document (\S\ref{sec:tbd}). The ROB treats \texttt{rd\_dbuf\_base} as an opaque base pointer whose only defined operation is \texttt{+ PACKET\_NUMBER}.
\end{itemize}

\subsection{Read Data Buffer (cross-reference)}
\begin{itemize}[noitemsep]
  \item The pooled Read Data Buffer (RD\_SRAM), its independently managed 128-entry watermark, and its occupancy/backpressure logic are \textbf{out of scope for this file}.  It is a separate structure the read ROB addresses through \texttt{rd\_dbuf\_base}, not a substructure of Metadata RAM. Data-buffer capacity and watermark are independent of \texttt{ROB\_DEPTH}/Metadata RAM capacity.
  \item Exact Data Buffer interface signal names, allocation representation, and timing are \textbf{TBD} in the MC/Data-Buffer interface contract --- see \S\ref{sec:tbd}.
\end{itemize}

\subsection{Read Serializer}
\begin{itemize}[noitemsep]
  \item Operates only while the Retiring Register (\S\ref{sec:retire}) holds a read transaction.
  \item Serializes the completed data at \texttt{rd\_dbuf\_base + PACKET\_NUMBER} for \texttt{PACKET\_NUMBER = 0 .. exp\_pkts-1} in packet order, tagged with the AXID selected from the Slot FIFO context.  Each 64-B MC packet is one AXI data beat; the within-beat handshake to the merge path follows the MC/Data-Buffer interface contract.
  \item \texttt{merge\_ready} accepts the AXI beat currently presented (one per 64-B MC packet).  The contiguous \texttt{rd\_dbuf\_base} range is released as a single unit once the final beat has been accepted.
  \item \texttt{rob\_last} is asserted on the final beat, i.e.\ when \texttt{PACKET\_NUMBER == exp\_pkts - 1} is accepted.
  \item Once the final beat is accepted and the \texttt{rd\_dbuf\_base} range is released, Retirement Logic pops the Slot FIFO and frees \texttt{ROB\_INDEX} (\S\ref{sec:retire}).
\end{itemize}

\subsection*{Read-Instance I/O (in addition to generic I/O in \S\ref{sec:io})}
\begin{longtable}{p{3cm}p{1cm}p{3cm}p{7cm}}
\toprule
\textbf{Signal} & \textbf{Width} & \textbf{Dir / Peer} & \textbf{Description} \\
\midrule
\texttt{data} & 512 & in, MC Core & Read response data, one 64-B unit per completed 64-B MC packet \\
\texttt{data} & 512 & out, Merge Logic & In-order read response data unit (one 64-B AXI R-channel beat) \\
\texttt{status[3:0]} & 4   & out, Merge Logic & Status for emitted beat \\
\texttt{rob\_last}   & 1   & out, Merge Logic & Asserted on the final AXI beat (\texttt{PACKET\_NUMBER == exp\_pkts - 1}) \\
\texttt{merge\_ready} & 1  & in, Merge Logic & Accept for the AXI beat currently emitted by the Read Serializer \\
\bottomrule
\caption{Read-instance additional signals}
\end{longtable}

\section{Write-Instance Extension}
\label{sec:write}

State and logic in this section exist only in the write-path ROB instance.

\subsection{Metadata RAM Extension}
\begin{itemize}[noitemsep]
  \item \texttt{resp\_status[1:0]}: aggregate \texttt{BRESP} value for the transaction. Initialized to \texttt{OKAY} at allocation; updated by Completion Input as each packet's status arrives, worst-of semantics (priority ordering \texttt{DECERR} $>$ \texttt{SLVERR} $>$ \texttt{OKAY}, encoding frozen in \S\ref{sec:cifmc}).
  \item \texttt{wr\_line\_base[25:0]}: start 64\,B line number (\texttt{ADDR\_W - 6}) of the write's footprint.
  \item \texttt{wr\_line\_count[6:0]}: contiguous 64\,B line span of the write, in \texttt{[1,64]}.
\end{itemize}
Both footprint fields are written by the allocation path from \texttt{\{AWADDR, AWLEN, AWSIZE, AWBURST\}} (\S\ref{sec:samelinehaz}), are never written by Completion Input, are read only by the read-admission same-line hazard check, and remain valid for the whole allocated lifetime of the entry --- released with it (\S\ref{sec:release}).

\subsection{BRESP Generation}
\begin{itemize}[noitemsep]
  \item Operates only while the Retiring Register (\S\ref{sec:retire}) holds a write transaction.
  \item Emits a single \texttt{bresp = resp\_status} to the Response Dispatcher, tagged with the AXID currently being serviced, and holds it until accepted downstream (\texttt{bresp\_ready}, \S\ref{sec:io}).
  \item No data movement, no per-packet serialization --- one response per transaction regardless of \texttt{exp\_pkts}.
  \item Once \texttt{bresp\_ready} accepts the response, Retirement Logic pops the Slot FIFO and frees \texttt{ROB\_INDEX} (\S\ref{sec:retire}) --- not before.
\end{itemize}

\subsection*{Write-Instance I/O (in addition to generic I/O in \S\ref{sec:io})}
\begin{longtable}{p{3cm}p{1cm}p{3cm}p{7cm}}
\toprule
\textbf{Signal} & \textbf{Width} & \textbf{Dir / Peer} & \textbf{Description} \\
\midrule
\texttt{bresp[1:0]} & 2 & out, Response Dispatcher & Aggregate write response for the transaction \\
\texttt{bresp\_ready} & 1 & in, Response Dispatcher & Accept for \texttt{bresp}; gates Slot FIFO pop and \texttt{ROB\_INDEX} free in Retirement Logic \\
\bottomrule
\caption{Write-instance additional signals}
\end{longtable}

\section{Generic I/O}
\label{sec:io}

\begin{longtable}{p{3cm}p{1cm}p{3cm}p{7cm}}
\toprule
\textbf{Signal} & \textbf{Width} & \textbf{Dir / Peer} & \textbf{Description} \\
\midrule
\texttt{alloc\_req}       & 1  & in, Accept/Request logic & Pulse: allocate a new ROB slot \\
\texttt{AXID[5:0]}        & 6  & in, Accept/Request logic & AXID of the transaction being allocated; used only to select the Per-AXID Slot FIFO, not stored in Metadata RAM \\
\texttt{exp\_pkts[6:0]}   & 7  & in, Burst Splitter & Number of 64-B MC packets for this transaction, \texttt{1..64} \\
\texttt{ROB\_INDEX[7:0]}  & 8  & out, Accept/Request logic & Allocated slot index, returned for use in \texttt{GLOBAL\_TAG} generation \\
\texttt{MC\_RESP} & TBD & in, MC Core & MC completion/response unit; exact transport defined by the MC--CIF interface contract \\
\texttt{GLOBAL\_TAG[13:0]} & 14 & in, MC Core & \texttt{\{ROB\_INDEX[7:0], PACKET\_NUMBER[5:0]\}}, echoed unmodified \\
\texttt{status[3:0]}      & 4  & in, MC Core & Per-packet completion status; \texttt{0}/\texttt{1}/\texttt{2} $=$ OK/SLVERR/DECERR, codes \texttt{3..15} reserved $\rightarrow$ DECERR (\S\ref{sec:cifmc}) \\
\bottomrule
\caption{Generic signals, present on both instances}
\end{longtable}

\noindent Instance-specific downstream accept signals (\texttt{merge\_ready}, \texttt{bresp\_ready}) are listed with their respective instance's I/O table above, since only one instance's Retirement Logic consumes each.

\section{Metadata RAM Entry Definition}
\label{sec:entry}

\begin{center}
\begin{tabular}{lll}
\toprule
\textbf{Field} & \textbf{Width} & \textbf{Scope} \\
\midrule
\texttt{exp\_pkts[6:0]}           & 7 bits          & generic \\
\texttt{completion\_bitmap[63:0]} & 64 bits         & generic \\
\texttt{rd\_dbuf\_base}           & \texttt{RD\_SRAM\_ADDR\_W} bits (\textbf{TBD}) & read-only \\
\texttt{resp\_status[1:0]}        & 2 bits          & write-only \\
\texttt{wr\_line\_base[25:0]}     & 26 bits         & write-only \\
\texttt{wr\_line\_count[6:0]}     & 7 bits          & write-only \\
\bottomrule
\end{tabular}
\end{center}

\noindent Total entry width: \(71 + \texttt{RD\_SRAM\_ADDR\_W}\) bits (read instance --- \texttt{exp\_pkts} 7 + \texttt{completion\_bitmap} 64 + \texttt{rd\_dbuf\_base}), 106 bits (write instance --- 71 generic + 2-bit \texttt{resp\_status} + the 33-bit 64\,B line footprint, \S\ref{sec:samelinehaz}). AXID is available from Slot FIFO context at retirement and is not stored in the entry. The write entry stores the transaction's 64\,B line footprint; \texttt{rd\_dbuf\_base} is the read entry's only address-like field and is an opaque RD\_SRAM base. No "retiring" state is stored per entry either --- that state lives entirely in the single Retiring Register described in \S\ref{sec:retire}, since at most one transaction per instance is ever draining at a time.

\section{Allocation Path}
\label{sec:alloc}

\begin{enumerate}[noitemsep]
  \item Accept/Request logic sends \texttt{alloc\_req} with \texttt{AXID}, following AXI request acceptance upstream of this file. On the read path, \texttt{alloc\_req} is withheld while the AR Metadata FIFO head is held by the same-line hazard check (\S\ref{sec:samelinehaz}); no read entry is allocated during the hold.
  \item Allocator pops \texttt{ROB\_INDEX} from the Free List (priority encoder on the bitmask), clears the bit. \texttt{ROB\_INDEX} is now reserved but not yet visible to Retirement Logic.
  \item Once Burst Splitter Stage 1 computes \texttt{exp\_pkts} (the number of 64-B MC packets, \texttt{1..64}) for the transaction, Metadata RAM is written at \texttt{ROB\_INDEX}: \texttt{exp\_pkts} set, \texttt{completion\_bitmap} cleared to \texttt{0}, (read instance) \texttt{rd\_dbuf\_base} set to the base of a contiguous \texttt{exp\_pkts}-line RD\_SRAM allocation, and (write instance) \texttt{resp\_status} initialized to \texttt{OKAY} with \texttt{wr\_line\_base} and \texttt{wr\_line\_count} written from \texttt{\{AWADDR, AWLEN, AWSIZE, AWBURST\}} (\S\ref{sec:samelinehaz}).
  \item Only \emph{after} this Metadata RAM initialization write completes is \texttt{ROB\_INDEX} pushed onto the requesting AXID's Slot FIFO. This ordering is deliberate: the Slot FIFO push and the Metadata RAM initialization write are both gated on the same Burst Splitter Stage 1 completion, so an entry can never become visible to Retirement Logic with a stale or uninitialized \texttt{exp\_pkts}/\texttt{completion\_bitmap}.
\end{enumerate}

\section{Packet Issue Path}
\label{sec:issue}

\begin{enumerate}[noitemsep]
  \item For each 64-B MC packet \texttt{PACKET\_NUMBER = 0..exp\_pkts-1} emitted by the Burst Splitter:
  \item Packet Formation attaches \texttt{GLOBAL\_TAG = \{ROB\_INDEX[7:0], PACKET\_NUMBER[5:0]\}} to the packet. No per-packet Read Data Buffer handle is recorded --- the contiguous \texttt{rd\_dbuf\_base} range was allocated at \S\ref{sec:alloc}, and packet \texttt{PACKET\_NUMBER}'s data will land at \texttt{rd\_dbuf\_base + PACKET\_NUMBER}.
  \item Packet is issued to MC Core via the Packet FIFO.
\end{enumerate}

\section{Completion Path}
\label{sec:completion}

\begin{enumerate}[noitemsep]
  \item The MC response unit arrives with \texttt{GLOBAL\_TAG},
      \texttt{status}, and, for reads, response data according to
      the MC--CIF interface contract.
  \item Completion Input decodes \texttt{GLOBAL\_TAG} into \texttt{ROB\_INDEX} and \texttt{PACKET\_NUMBER}.
  \item (read instance) the returned 64-B data unit is written to \texttt{rd\_dbuf\_base + PACKET\_NUMBER} in RD\_SRAM.
  \item (write instance) \texttt{resp\_status} is updated with the worst-of comparison against the incoming \texttt{status}.
  \item \texttt{completion\_bitmap[PACKET\_NUMBER]} is set at \texttt{ROB\_INDEX}.
  \item Completion acceptance only marks the corresponding packet
      complete; it never triggers a Slot FIFO pop or a Free List return.
      Retirement remains responsible for releasing the ROB entry.
\end{enumerate}

\section{Retirement Path}

\begin{enumerate}[noitemsep]
  \item Each cycle, while the Retiring Register is idle, Retirement Logic reads the Metadata RAM entry at every AXID's Slot FIFO head.
  \item Priority encoder selects the lowest AXID whose head entry has \texttt{completion\_bitmap} bits \texttt{[exp\_pkts-1:0]} all set; its \texttt{ROB\_INDEX} is loaded into the Retiring Register. \textbf{The Slot FIFO entry is not popped yet.}
  \item (read instance) Read Serializer drains \texttt{rd\_dbuf\_base + 0 .. rd\_dbuf\_base + exp\_pkts-1} in packet order, one AXI beat per 64-B MC packet.
  \item (write instance) BRESP Generation emits \texttt{bresp} and holds it until \texttt{bresp\_ready} accepts it.
  \item Only once the transaction's full AXI response has been accepted (final read beat, or the single write response) does Retirement Logic pop that AXID's Slot FIFO and return \texttt{ROB\_INDEX} to the Free List. The Retiring Register returns to idle.
\end{enumerate}

\section{Slot Release}
\label{sec:release}

\begin{enumerate}[noitemsep]
  \item Triggered only by full AXI response acceptance (\S\ref{sec:retire}) --- never by \texttt{completion\_bitmap} reaching its full mask on its own.
  \item \texttt{ROB\_INDEX} is returned to the Free List (\texttt{free\_req}), bit set back to \texttt{1 = free}.
  \item (write instance) the \texttt{free\_req} pulse re-arms the same-line hazard recheck for any AR Metadata FIFO head currently held (\S\ref{sec:samelinehaz}); the freed entry's \texttt{wr\_line\_base} and \texttt{wr\_line\_count} cease to be occupied in the same cycle.
  \item (read instance) the contiguous \texttt{rd\_dbuf\_base} range (\texttt{exp\_pkts} lines) is returned to RD\_SRAM as a single unit, once the final AXI beat has been accepted.
  \item Retirement needs only the Slot FIFO pop and Free List return.
\end{enumerate}

\section{Backpressure}
\label{sec:backpressure}

\begin{center}
\begin{tabularx}{\textwidth}{>{\ttfamily}l X}
\toprule
\textbf{Signal} & \textbf{Description} \\
\midrule
Slot FIFO full (per AXID) & Blocks new allocation for that AXID. A completed-but-still-draining head entry keeps the FIFO occupied until fully retired, so this condition can now persist slightly longer than "completion" alone would suggest --- expected behavior, not a fault. \\
FREE\_LIST\_EMPTY         & Blocks new allocation pool-wide. It means all 256 transaction slots are allocated and all Slot FIFO capacity is collectively consumed. A single full AXID FIFO does not imply Free List exhaustion. \\
AR Metadata FIFO head held & Read path only. The head read overlaps an occupied \texttt{ROB\_WR} 64\,B line footprint (\S\ref{sec:samelinehaz}); read admission on that port is held until every overlapping write has retired. No write-path equivalent. \\
\bottomrule
\end{tabularx}
\end{center}

\medskip
\noindent \texttt{ROB\_DEPTH = 256} carries no watermark of its own. Read-path issuance backpressure is driven entirely by the pooled Read Data Buffer's 128-entry watermark, a separate structure not detailed in this file.

\section{Key Parameters}
\label{sec:params}

\begin{center}
\begin{tabular}{ll}
\toprule
\textbf{Parameter} & \textbf{Value} \\
\midrule
Metadata RAM depth (\texttt{ROB\_DEPTH}) & 256 entries (\texttt{ROB\_INDEX} is 8 bits) --- frozen, identical for both instances \\
\texttt{PACKET\_NUMBER} width & 6 bits (\texttt{MAX\_PKTS = 64}) --- frozen \\
\texttt{GLOBAL\_TAG} width    & 14 bits --- frozen \\
\texttt{exp\_pkts} width      & 7 bits (packet count \texttt{1..64}) --- frozen \\
\texttt{completion\_bitmap} width & 64 bits --- frozen \\
Entry width, read instance    & \(71 + \texttt{RD\_SRAM\_ADDR\_W}\) bits (\texttt{RD\_SRAM\_ADDR\_W} \textbf{TBD}) \\
Entry width, write instance   & 106 bits (71 generic + 2-bit \texttt{resp\_status} + 33-bit 64\,B line footprint, \S\ref{sec:samelinehaz}) \\
Free list width          & 256 bits, one independent Free List per instance \\
Per-AXID Slot FIFO        & 64 FIFOs $\times$ depth 4 --- frozen, identical for both instances \\
Retiring Register         & 1 per instance: \texttt{ROB\_INDEX[7:0]} + 6-bit read packet pointer + within-beat progress state \\
Max transactions draining & 1 per instance at a time; a single read transaction's drain spans one AXI beat per 64-B MC packet, up to \texttt{exp\_pkts} $\le$ 64 beats \\
\bottomrule
\end{tabular}
\end{center}

\noindent \(64 \times 4 = 256\), exactly matching \texttt{ROB\_DEPTH}. Thus \texttt{FREE\_LIST\_EMPTY} means all per-AXID FIFO capacities are collectively occupied; it remains distinct from the condition that any one FIFO is full.

\section{Unresolved / TBD}
\label{sec:tbd}

Every item below requires an external decision or a separate component's specification --- none can be derived from the frozen ROB architecture itself.

\begin{itemize}[noitemsep]
  \item \textbf{MC error-code assignment (codes \texttt{3..15}).} The interface \texttt{status[3:0]} field and its AXI mapping are frozen (\S\ref{sec:cifmc}): \texttt{0}/\texttt{1}/\texttt{2} $=$ OK/SLVERR/DECERR, and any received value in \texttt{3..15} is treated by CIF as \texttt{DECERR}. Which specific fault classes MC Core assigns to codes \texttt{3..15}, if any, is an MC-owned detail and does not affect CIF RTL.
  \item \textbf{Pooled Read Data Buffer / MC response contract.} Signal names, allocation representation, and timing for allocating, filling, serializing, and freeing the read data range are defined by a separate MC/Data-Buffer interface specification. This file assumes only that the contiguous \texttt{rd\_dbuf\_base} range is valid from allocation until the contract-defined release point, and that packet \texttt{N}'s returned 64-B data unit occupies \texttt{rd\_dbuf\_base + N}.
  \item \textbf{RD\_SRAM depth / \texttt{RD\_SRAM\_ADDR\_W}.} The depth of RD\_SRAM, and therefore the width of \texttt{rd\_dbuf\_base} and the read-instance entry, is not specified in any current document. The ROB carries \texttt{rd\_dbuf\_base} as an opaque base pointer; RD\_SRAM's depth, occupancy model, and watermark are out of scope for this file. The RD\_SRAM completion-side write port must be available whenever a \texttt{RD\_DATA} response is presented (\S\ref{sec:cifmc}, response-side flow control).
  \item \textbf{MC write-completion safety (interface requirement).} The property that a write reported complete is visible to a subsequent same-line read is stated as an interface requirement on MC Core (\S\ref{sec:cifmc-req}, item 2), \emph{not} asserted to be implemented. The current MC-core documentation additionally contains an unresolved contradiction about how long a write remains hazard-visible (its hazard-unit admission-station dwell versus its full buffer-to-DRAM lifetime); the MC team must resolve that before the requirement can be signed off. Until then, the conservative interpretation --- \texttt{WR\_ACK} at DRAM commit --- satisfies the requirement without any MC hazard-unit dependency.
\end{itemize}

\section{Timing Notes}

\begin{itemize}[noitemsep]
  \item Metadata RAM read/write ports: allocation write, completion write, and retirement read are independent accesses; simultaneous allocation and completion in the same cycle to different \texttt{ROB\_INDEX} values is safe, unchanged in spirit from the prior BRAM port arrangement.
  \item Priority encoder across 64 Slot FIFO heads in Retirement Logic remains the likely critical path; may need pipelining at high frequencies, as in the prior design. It only runs while the Retiring Register is idle.
  \item Simultaneous \texttt{alloc\_req} and \texttt{free\_req} on the Free List bitmask in the same cycle is safe --- they target different bits by construction (allocate clears the lowest free bit; free sets the bit of the just-retired, and therefore currently occupied, slot).
  \item \texttt{GLOBAL\_TAG} decode (\S\ref{sec:tag}) is a pure bit-slice and has no translation latency on the completion path.
  \item The Slot FIFO push at allocation is deliberately deferred until the Metadata RAM initialization write completes (\S\ref{sec:alloc}) --- this introduces a small, bounded latency between \texttt{alloc\_req} and Slot FIFO visibility, bounded by Burst Splitter Stage 1's latency to produce \texttt{exp\_pkts}.
  \item Draining a read transaction is explicitly multi-cycle through the Retiring Register.  Transfer granularity is one AXI beat per 64-B MC packet, so Retirement Logic and the Read Serializer must be read as a small state machine (idle / retiring), not a single combinational selection. The write instance is analogous but single-beat, gated by \texttt{bresp\_ready}.
\end{itemize}